\section{ผลการตรวจระบบและแนวทางประเมิน}
\subsection{วิธีตรวจและขอบเขตของหลักฐาน}
ตรวจโครงสร้างหน้าเว็บ สคีมาฐานข้อมูล บริการจัดการใบลา ระบบข้อมูลสนับสนุน และโค้ด ETL กับบริการรายงาน แล้วจับคู่กับวัตถุประสงค์ในข้อเสนอโครงงาน นำผลตรวจมาจัดข้อกำหนดตามบทบาท นอกจากนี้อ่าน metadata และ SHOW CREATE TABLE ของฐานใบลาและ warehouse บน Railway เมื่อ 18 ต.ค. 2569 โดยไม่อ่านแถวข้อมูลธุรกรรม การตรวจนี้ไม่ได้ทดสอบการทำงานของแอปพลิเคชัน
\subsection{ข้อกำหนดการทำงานตามบทบาท}
ตารางที่~\ref{tab:functional-role} กำหนด functional requirements (FR) ตามผู้ใช้งานและขอบเขตข้อมูล ใช้เป็นฐานออกแบบกรณีทดสอบ โดยการพบโค้ดรองรับไม่ได้หมายถึงผ่านการทดสอบแล้ว ผู้ใช้หนึ่งคนอาจมีหลายบทบาท แต่ทุกคำขอต้องตรวจตัวตนและความสัมพันธ์กับข้อมูลฝั่งบริการ
\begin{table}[ht]
\caption{ข้อกำหนดการทำงานและบทบาท}\label{tab:functional-role}
\centering
\begin{tabularx}{\columnwidth}{|p{23mm}|X|}
\hline
รหัส / บทบาท & ข้อกำหนดการทำงาน \\ \hline
FR01 นิสิต & ยื่นใบลา เลือกวัน ประเภท เหตุผล วิชาและตอนเรียน; ตรวจข้อมูลจำเป็น \\ \hline
FR02 นิสิต & ติดตามประวัติและสถานะของตน จัดการหลักฐานตามสถานะ และส่งออกใบลา \\ \hline
FR03 ที่ปรึกษา & พิจารณาใบลาที่อยู่ในขั้นของตน และดูรายงานนิสิตในความดูแล \\ \hline
FR04 หัวหน้าภาค / ผู้บริหาร & พิจารณาตามลำดับและดูรายงานตามขอบเขตหน่วยงานที่มีสิทธิ์ \\ \hline
FR05 ผู้สอน & อ่านและรับทราบใบลาที่อนุมัติแล้วของวิชาที่เกี่ยวข้อง พร้อมดูรายงานรายวิชา \\ \hline
FR06 เจ้าหน้าที่ & อ่านและติดตามข้อมูลใบลาและรายงานตามสิทธิ์ที่ระบบกำหนด \\ \hline
FR07 ผู้ดูแลข้อมูล & จัดการวิชา ตอนเรียน ผู้สอนและที่ปรึกษา; นำเข้า Excel แบบเพิ่มหรือแทนที่ \\ \hline
FR08 ระบบ & คุมลำดับอนุมัติ ป้องกันตัดสินซ้ำ บันทึกคำตัดสินและแหล่งที่มา รวมกฎอัตโนมัติ \\ \hline
FR09 ระบบรายงาน & โหลดข้อมูลแบบ snapshot แสดงเวลาสำเร็จล่าสุด กรองตามบทบาท และนับใบลาไม่ซ้ำ \\ \hline
\end{tabularx}
\end{table}
รายละเอียดเกณฑ์ยอมรับราย FR และแผนบันทึกหลักฐานอยู่ในภาคผนวกท้ายเอกสาร เพื่อคงบทความหลักไม่เกิน 6 หน้า

ผลการตรวจสนับสนุนว่าระบบมีองค์ประกอบสำหรับกระบวนการที่กำหนด แต่ยังไม่เพียงพอที่จะสรุปว่าการทำงานทุกกรณีถูกต้อง หรือผู้ใช้พึงพอใจในระดับใด ภาพหน้าจอจากระบบจริงและหลักฐานการทดสอบควรจัดเก็บเพิ่มเติมก่อนส่งบทความฉบับสมบูรณ์
\subsection{แผนทดสอบการทำงาน}
ควรทดสอบการยื่นคำขอที่ข้อมูลครบและไม่ครบ การเข้าถึงคำขอของนิสิตอื่น การพิจารณาผิดลำดับ การส่งผลตัดสินซ้ำ และการตัดสินพร้อมกัน รวมถึงการแนบหลักฐานก่อนและหลังมีการพิจารณา การรับทราบของผู้สอน และการกรองรายงานตามบทบาท โดยใช้ข้อมูลทดสอบที่แยกจากข้อมูลจริง
ควรทดสอบ ระบบข้อมูลสนับสนุน เพิ่มทั้งการอ่านวิชาและตอนเรียน การนำเข้าแบบเพิ่มและแทนที่ การแก้ที่ปรึกษา ข้อมูลซ้ำ และการเรียกบริการไม่ได้ จากนั้นตรวจผลต่อแบบฟอร์มและขอบเขตรายงาน โดยไม่เปลี่ยนข้อมูลประวัติใบลาที่ส่งแล้ว การตรวจ ETL ควรครอบคลุมการรันซ้ำ การเปลี่ยนข้อมูลลูก การ rollback และ reconcile รวมถึงเวลาการโหลดสำเร็จล่าสุด
การตรวจรายงานควรใช้ชุดข้อมูลที่ทราบคำตอบ เช่น ใบลาหนึ่งใบเกี่ยวข้องกับหลายวิชา แล้วตรวจว่าจำนวนใบลาไม่เพิ่มตามจำนวนวิชา พร้อมทดสอบกรณีไม่มีรายการและกรณีข้อมูลมิติไม่ครบ บันทึกผลคาดหวัง ผลจริง และสภาพแวดล้อมของแต่ละกรณีเพื่อให้ตรวจซ้ำได้
\subsection{เกณฑ์ยอมรับสำหรับการทดสอบร่วมกัน}
การทดสอบร่วมกันควรติดตามเส้นทางเดียวตั้งแต่จัดเตรียมข้อมูลอ้างอิง ยื่นคำขอ พิจารณาจนสิ้นสุด โหลดข้อมูล แล้วอ่านรายงาน เพื่อหาความคลาดเคลื่อนที่การตรวจแต่ละส่วนแยกกันอาจไม่พบ ตารางที่~\ref{tab:acceptance} เป็นกรณีที่เสนอและผลคาดหวัง ยังไม่มีสถานะผ่านหรือไม่ผ่านจากการทดลองในบทความนี้
\begin{table}[ht]
\caption{กรณีทดสอบร่วมกันและผลคาดหวัง}\label{tab:acceptance}
\centering
\begin{tabularx}{\columnwidth}{|p{25mm}|X|}
\hline
กรณีที่เสนอ & ผลคาดหวังที่ต้องตรวจ \\ \hline
หนึ่งใบหลายวิชา & จำนวนใบและวันระดับคำขอไม่เพิ่มตามวิชา \\ \hline
ผู้สอนร่วม & ผู้เกี่ยวข้องเข้าถึงตามสิทธิ์และไม่เพิ่มยอดใบลาซ้ำ \\ \hline
เปลี่ยนที่ปรึกษา & สิทธิ์ปัจจุบันเปลี่ยนตามต้นทาง; ประวัติเดิมยังมีบริบท \\ \hline
โหลดซ้ำชุดเดิม & คีย์และยอดข้อมูลเดิมไม่เพิ่มซ้ำ \\ \hline
โหลดล้มเหลว & ไม่เผยแพร่ข้อมูลบางส่วนเป็น batch สำเร็จ \\ \hline
เทอมไม่มีใบลา & มีค่าเป็นศูนย์เมื่อเป็นเทอมที่เลือก \\ \hline
\end{tabularx}
\end{table}
ชุดข้อมูลทดสอบควรกำหนดคำตอบล่วงหน้าและแยกกรณีสิทธิ์ออกจากกรณีรวมยอด ตัวอย่างข้อมูลสมมติสองใบจำนวน 3 วันเท่ากันช่วยตรวจข้อผิดพลาดจากการตัดซ้ำด้วยค่าจำนวนวัน ส่วนการเพิ่มผู้สอนร่วมช่วยตรวจผลคูณของแถวเมื่อเชื่อมตาราง ต้องบันทึกยอดที่ระดับคำขอและระดับวิชาแยกกันเพื่อระบุจุดที่คลาดเคลื่อน
